\documentclass[11pt,a4paper]{article}

\usepackage[margin=2.5cm]{geometry}
\usepackage[T1]{fontenc}
\usepackage{lmodern}
\usepackage{amsmath}
\usepackage{booktabs}
\usepackage{tabularx}
\usepackage{enumitem}
\usepackage{xcolor}
\usepackage{listings}
\usepackage{hyperref}

\hypersetup{colorlinks=true, linkcolor=blue!50!black, urlcolor=blue!50!black}
\lstset{basicstyle=\ttfamily\small, columns=fullflexible, breaklines=true, frame=single,
  framerule=0.3pt, rulecolor=\color{black!30}, backgroundcolor=\color{black!3}, showstringspaces=false}
\newcommand{\code}[1]{\texttt{#1}}

\title{CPU--GPU Co-Scheduling Module for CloudSim 7G\\[4pt]
\large Phase 5 Report: Energy and CO$_2$, Comparisons and the Decision Report}
\author{Chirag \and Hamiz}
\date{26 November 2026}

\begin{document}
\maketitle

\begin{abstract}
Phases 0--4 built a simulator that reproduces Synergy's simulator job for job. Phase~5 turns it into the
decision tool the architecture describes. Each run now measures cluster utilisation and, when power inputs are
given, IT and facility energy, CO$_2$, energy per job and per training sample, and the energy drawn by GPUs that
are idle or waiting for data; these integrals are exact, because power only changes when the schedule does.
A seeded Poisson workload generator makes seeds meaningful, sweeps answer what-if questions such as
``24 or 32 cores per server?'', and a \code{ComparisonRunner} runs every candidate $\times$ seed $\times$ sweep
value. A \code{DecisionReport} then ranks candidates per metric and by the researcher's weights, marks the
trade-offs, says whether each lead is clear or within the noise between seeds, flags leaders that change across
sweep values, lists every assumption, and writes everything as LaTeX tables and CSV. Power numbers have no
defaults: every value must be written in the experiment file with its source. A Maven wrapper and a
\code{./cosched} script make the project run from a fresh clone with only Java installed. All 290 tests pass,
and the Synergy validation still matches exactly.
\end{abstract}

\tableofcontents

%======================================================================
\section{What was added}
%======================================================================

\begin{center}
\begin{tabularx}{\linewidth}{@{}lX@{}}
\toprule
Piece & What it does \\
\midrule
\code{PowerSettings}, \code{energy} section & Power inputs per server group and PUE / grid factor, all required with sources \\
\code{ClusterMeter} (L5) & The architecture's EnergyListener plus utilisation: exact integrals over the run, \code{timeline.csv} \\
\code{PoissonWorkload} (L3) & Seeded arrivals at a given load; jobs from Synergy's synthetic recipe or resampled from a trace \\
\code{workload.measure} & JCT statistics over a window of job ids (steady state), as Synergy measures \\
\code{Sweeps} & What-if axes: set any value in the YAML, several settings together if needed \\
\code{ComparisonRunner} (L5) & Validates every sweep point, then runs candidates $\times$ seeds $\times$ sweep values \\
\code{MetricCatalog}, \code{DecisionReport} (L5) & Per-metric ranks, weighted score, Pareto set, stability, LaTeX + CSV \\
\code{./mvnw}, \code{./cosched} & Build with only Java 21; one command runs an experiment and its report \\
\bottomrule
\end{tabularx}
\end{center}

The GPU layer gained one thing: \code{ServerStatus} now reports the CPU and GPU parts of a server's power as well
as the total. Scheduling is unchanged; the broker only gained a listener that is called after every change to
what runs where.

%======================================================================
\section{Energy and CO$_2$}
%======================================================================

\subsection{Model}
As in the architecture document,
\begin{equation*}
E_{\text{IT}} = \sum_{\text{servers}} \int \Big( P_{\text{CPU}}(u_{\text{CPU}}) + \sum_{\text{GPUs}} P_{\text{GPU}}(u_{\text{GPU}})
  + P_{\text{base}} \Big)\, dt, \qquad E_{\text{facility}} = E_{\text{IT}} \times \text{PUE}, \qquad
\text{CO}_2 = E_{\text{facility}} \times \text{grid factor}.
\end{equation*}
$u_{\text{CPU}}$ is the share of cores allocated (allocated cores are assumed busy). A GPU's busy fraction is its
job's speedup divided by its speedup at the ideal allocation, where the ideal is capped at one server as Synergy
caps it. $P_{\text{CPU}}$ is linear from idle to maximum, or one of CloudSim's SPECpower tables;
$P_{\text{GPU}}$ is linear; idle servers stay on.

\subsection{Why the integral is exact}
Allocations, busy fractions and therefore power change only when a plan is applied or a job ends. The broker
calls the meter at exactly those moments, and the meter multiplies the power held since the previous call by
the time elapsed. There is no sampling interval to choose and no approximation; the run's window is
$[0, \text{makespan}]$.

\subsection{Where GPU energy goes}
A GPU busy a fraction $u$ of the time draws idle power while waiting and maximum power while computing, so
$(1-u)\,P_{\text{idle}}$ of its power is spent waiting. The meter splits that waiting energy into
\emph{idle} (no job holds the GPU) and \emph{stalled} (a job holds it but is starved of CPU or memory, so the GPU
waits for data). Stalled energy is what CPU-aware scheduling can remove.

\subsection{No defaults, every number with a source}
Energy is modelled only when the YAML has an \code{energy} section, and then every server group must have a
\code{power} section. Each number must be written, and a \code{source} field must say where the numbers come from:
\begin{lstlisting}
power: { cpu_model: linear, cpu_idle_w: 100, cpu_max_w: 300, gpu_idle_w: 50, gpu_max_w: 300, base_w: 200,
         source: "PLACEHOLDER values, not measured ..." }
energy: { pue: 1.54, pue_source: "...", grid_g_co2_per_kwh: 727, grid_source: "..." }
\end{lstlisting}
A missing number, a missing source, power without an \code{energy} section, or idle/max watts given together with
a SPECpower table are errors before the run. The sources appear in the manifest and in the report's assumptions,
so a placeholder can never pass for a measurement. The demo preset uses clearly labelled placeholders; PUE 1.54
and 727\,gCO$_2$/kWh are the values the architecture document cites (Uptime Institute 2025; CEA v20.0).

\subsection{Metrics}
IT and facility energy (kWh), CO$_2$ (kg), mean IT power, energy per finished job, energy per training sample
(iterations $\times$ the profile's \code{batch\_size}), and the shares of IT energy drawn by CPUs, by GPUs, by idle
GPUs and by stalled GPUs. Utilisation, measured in every run: time-averaged GPUs allocated, GPUs busy and cores
allocated.

\subsection{Check by hand}
One 8-GPU, 24-core server; CPU 100--300\,W, GPU 50--250\,W, base 100\,W. A transformer job holds all 8 GPUs and 24
cores for 100\,s with fully busy GPUs (2400\,W). Then a ResNet-18 job holds 4 GPUs and 12 cores for 100\,s at the
fixed share. Its ideal (9 cores per GPU, 36 cores) is more than a server has, so it is capped at 6 cores per GPU,
speedup 1.75: its GPUs are busy $1/1.75$ of the time. The power is
$200 + 4(50 + 200/1.75) + 4 \times 50 + 100 = 1157.1$\,W. The test checks IT and facility energy, CO$_2$, energy
per job and per sample, the idle, stalled and CPU shares, and the three utilisations against this calculation.
Writing the test first with an uncapped busy fraction ($1/2.4$) gave a 3.6\% mismatch; the code was right,
because it applies Synergy's cap, and the test was corrected.

%======================================================================
\section{A seeded workload}
%======================================================================

Replaying a fixed trace gives the same result for every seed, so seeds had no effect. The \code{poisson}
generator draws arrivals as Synergy does (\code{helpers/utils.py}): each gap is
$\lceil -\ln(1-U)/\lambda \times 3600 \rceil$ whole seconds at $\lambda$ jobs per hour. What each job is comes
from one of two sources:
\begin{itemize}
  \item \code{synthetic}: Synergy's synthetic recipe. Run time at the fixed share is $60 \times 10^{U(1.5,3)}$\,s
        with probability 0.8, else $60 \times 10^{U(3,4)}$\,s; GPUs are 1, or Synergy's multi-GPU mix (1: 70\%,
        2: 10\%, 4: 15\%, 8: 5\%). Models come from the labeller.
  \item \code{resample}: each job copies the GPUs, work and model of a random row of a job CSV. Resampling
        Synergy's reference job list gives the Philly-derived durations its paper uses.
\end{itemize}
The load and the job count are required: they define the experiment. For a given seed every candidate runs on
the same jobs (common random numbers), so differences between candidates are paired. The jobs are written to
each run's \code{workload.csv}, so any run can be replayed exactly. \code{workload.measure} restricts the JCT
statistics to a window of job ids, the steady state, as Synergy measures jobs 4000--4999; energy and utilisation
always cover the whole run.

%======================================================================
\section{Sweeps and the comparison runner}
%======================================================================

A sweep sets one value in the YAML tree to each of its values in turn; several sweeps multiply. Because it
edits the YAML, anything written there can be swept: a cluster dimension, a policy parameter, the load. Some
settings must change together. More cores per server also need profiles normalised to that server's fixed
share (Synergy's \code{models-4} for 32 cores), so one sweep can set several paths:
\begin{lstlisting}
sweeps:
  - name: cpus
    paths: [cluster.server_groups.0.cpus, models.profiles_dir]
    values: [[24, ../profiles/synergy-v100/models], [32, ../profiles/synergy-v100/models-4]]
\end{lstlisting}
\code{ComparisonRunner} expands the sweeps and validates every combination before the first run: a wrong path,
a misspelt key or a bad value stops the study at once rather than after an hour of simulation. It then runs
every candidate $\times$ seed at every point, under \code{out/<experiment>/<point>/<candidate>/seed-<seed>/}, and
writes the report. Runs are sequential, since CloudSim keeps its simulation state in static fields.

%======================================================================
\section{The decision report}
%======================================================================

The architecture asks for a report that ``ranks policies per metric, shows trade-offs, and flags rankings that
change across seeds or load levels instead of naming a single winner''. Each piece:

\paragraph{Metric catalogue.} Every metric has a unit, a print format, a meaning, and a direction (lower better,
higher better, or none). Two metrics that look rankable are shown only as context: \emph{GPUs allocated} and
\emph{idle-GPU energy}. In the first demo run the GPU-proportional baseline ``won'' both, only because its jobs
run slowly and hold their GPUs longer. Ranking them would reward slowness.

\paragraph{Per-metric ranks} by mean over seeds, with the 95\% confidence interval of each mean.

\paragraph{Weighted score.} For each objective, a candidate's \emph{regret} is how much worse its mean is than
the best candidate's, relative to the best: $(x - x_{\text{best}})/|x_{\text{best}}|$ for lower-is-better metrics
(when the best is 0, the gap is divided by the range). The score is the weighted mean regret: 0\% means best on
every objective, and 10\% means 10\% worse than the best on weighted average. Unlike a rank, this keeps the size of
a gap. Without objectives in the YAML there is no weighted ranking: the tool does not choose weights for the
researcher.

\paragraph{Trade-offs.} A candidate is Pareto-optimal if no other candidate is at least as good on every
objective and better on one.

\paragraph{Stability across seeds.} For each metric: the leader by mean; the seeds on which it is also best; and
the gap to the runner-up as a paired 95\% $t$-interval over seeds. The verdict is \emph{clear} if the interval
excludes 0, \emph{same order on every seed, size uncertain} if the leader wins every seed but the gap varies too
much to pin down, and \emph{within seed noise} otherwise (a tie at this number of seeds).

\paragraph{Stability across sweep values.} The leader of every metric at every point, with metrics whose leader
changes marked: for those, the answer to the what-if question is ``it depends''.

\paragraph{Outputs.} \code{decision-report.tex} (standalone, compiles with pdflatex) with the candidates table,
per point the metrics, ranking and stability tables, the sweep table, the assumptions and the metric definitions.
Each table is also its own file under \code{tables/} for use in other documents. Every number is in
\code{runs.csv}, \code{summary.csv}, \code{ranking.csv}, \code{stability.csv} and \code{sweep\_stability.csv}.

%======================================================================
\section{Demonstration: \code{presets/decision-demo}}
%======================================================================

A small what-if study built from the same plug-ins as the replication. Four Synergy-shaped servers (8 V100,
500\,GB) with 24 or 32 cores; Poisson arrivals at 2.25 jobs per hour (Synergy's 9 per hour scaled to a quarter of
its cluster), 1500 jobs resampled from Synergy's reference job list, JCTs over jobs 500--1199; candidates FIFO with
the fixed share, Greedy and Tune, and SRTF with Tune; 5 seeds; objectives mean JCT (weight 2), p99 JCT (1) and
energy per job (1). The 40 runs take 25\,s.

\begin{center}
\begin{tabular}{@{}llrrrr@{}}
\toprule
Cores & Candidate & Mean JCT (h) & p99 JCT (h) & kWh per job & Score \\
\midrule
24 & fifo-fair   & $55.8 \pm 23.1$ & $182 \pm 32$ & $8.66$ & 136.9\% \\
24 & fifo-greedy & $24.3 \pm 13.0$ & $146 \pm 23$ & $8.25$ & 28.2\% \\
24 & fifo-tune   & $23.7 \pm 13.2$ & $146 \pm 23$ & $8.24$ & 26.3\% \\
24 & srtf-tune   & $15.5 \pm 3.1$  & $183 \pm 63$ & $8.24$ & 6.4\% \\
\midrule
32 & fifo-fair   & $39.0 \pm 19.3$ & $161 \pm 29$ & $8.55$ & 82.4\% \\
32 & fifo-greedy & $21.5 \pm 12.0$ & $144 \pm 21$ & $8.15$ & 20.6\% \\
32 & fifo-tune   & $21.5 \pm 11.9$ & $144 \pm 22$ & $8.15$ & 20.6\% \\
32 & srtf-tune   & $15.2 \pm 2.7$  & $184 \pm 63$ & $8.14$ & 7.0\% \\
\bottomrule
\end{tabular}
\end{center}

What the report says, and what it does not:
\begin{itemize}
  \item CPU-aware allocation matters most where CPUs are scarce: with 24 cores, Tune's mean JCT is 58\% lower
        than the fixed share's; with 32 cores the fixed share improves by 30\% and the gap narrows.
  \item SRTF+Tune has the lowest mean JCT on every seed, but the size of its lead varies (``same order on every
        seed, size uncertain''). It also has the worst p99, as shortest-first scheduling is expected to. The
        report keeps that trade-off visible: SRTF, FIFO+Tune and Greedy are all Pareto-optimal, and only the fixed
        share is dominated.
  \item Tune and Greedy are within seed noise of each other at this load.
  \item Energy per job differs by under 6\% between candidates, and its leader is within noise. With idle servers
        kept on, a run's energy is set mostly by its length, which the last arrival fixes. The energy columns are
        context here, not a differentiator. The power numbers are placeholders, and the report says so.
  \item The p99.9 leader changes from FIFO+Tune (24 cores) to Greedy (32 cores); the sweep table flags it.
\end{itemize}

%======================================================================
\section{Running it}
%======================================================================

\begin{lstlisting}
git clone https://github.com/hMiZ-STD/Stacked-cpu-gpu-cloudsim-scheduler
cd Stacked-cpu-gpu-cloudsim-scheduler
./mvnw install                                       # Java 21 only; Maven 3.9.9 is downloaded
./cosched presets/decision-demo/decision-demo.yaml   # report: out/decision-demo/report/
\end{lstlisting}
\code{./cosched} rebuilds when sources have changed, builds the classpath once, and prints a short summary;
a mistake in the YAML gives a one-line error instead of a stack trace. The format of every new YAML section and
output file is in \code{docs/formats.md}.

%======================================================================
\section{Tests}
%======================================================================

\begin{center}
\begin{tabularx}{\linewidth}{@{}lX@{}}
\toprule
Test & Checks \\
\midrule
\code{EnergyTest} (4) & The hand calculation above; no energy metrics without an \code{energy} section, utilisation
  still measured; each missing number, missing source and invalid combination rejected; a SPECpower CPU table
  (IBM x3550, 222\,W at full load) integrates correctly. \\
\code{PoissonWorkloadTest} (6) & Same seed same jobs, other seed other jobs; mean gap 400\,s at 9 jobs/hour, whole
  seconds; duration bounds, 20\% long jobs, the GPU mix; resampling copies rows; a written workload reads back
  exactly; load and job count required, unused parameters rejected. \\
\code{SweepsTest} (4) & Nested keys, list positions and \code{*}; bad paths; a misspelt key caught by validation;
  product order and labels; several paths together; the original tree untouched. \\
\code{DecisionReportTest} (9) & Mean, sd and CI; ranks; regret scores and the Pareto set worked by hand; all
  three verdicts; one seed; missing runs; leader change across sweep values; every file written, balanced braces,
  LaTeX escaping; unknown or unrankable objectives rejected before running. \\
\code{ComparisonRunnerTest} (3) & $2 \times 2 \times 2$ runs with every file; one workload per seed shared by all
  candidates; the sweep changes the load; measured window; a bad sweep fails before any run; the single-run entry
  point refuses sweeps. \\
\bottomrule
\end{tabularx}
\end{center}

\begin{center}
\begin{tabular}{@{}lr@{}}
\toprule
Module & Tests (all passing) \\
\midrule
\code{cloudsim} (unchanged) & 167 \\
\code{cloudsim-examples} & 13 \\
\code{cloudsim-gpu} & 16 \\
\code{cloudsim-cosched} & 94 (1 skipped: the opt-in Synergy suite) \\
\midrule
Total & 290 \\
\bottomrule
\end{tabular}
\end{center}

The step-0 Synergy validation runs in every build and still matches exactly. The opt-in six-configuration subset
of the Phase~4 suite was rerun after these changes and passes (every measured job exact), so the new hooks
changed no scheduling decision.

%======================================================================
\section{Scope decisions and limits}
%======================================================================

\paragraph{Teams, quotas and fairness are out of this phase.} The configuration accepts teams and quotas and the
job CSV has a team column, but no policy uses them yet. Three reasons. It is a separate research question:
this project is about how many cores each job gets, while fair sharing between teams decides which team's jobs run,
and other work studies that (Themis, Gandiva-fair). It could not be validated: Synergy, the reference, has no
teams, and every result so far rests on an exact match. And it needs no redesign: a fair-share ordering is one new
policy class and one registry line. It is planned as an extension.

\paragraph{DRF and Tetris} (Synergy's Fig.~13) are the next policies. DRF is also the test of the
architecture's claim that a policy is one class plus one registry line with no other change, so it is assigned
to Hamiz, who did not write the pipeline, and will be validated against Synergy's simulator like the others.

\paragraph{Limits, stated in every report.} The power numbers in the presets are placeholders. A straight line to
rated maximum overestimates GPU power. The GPU busy fraction from throughput is our assumption. Idle servers stay
on, and grid carbon is constant. No published paper exists to validate the energy numbers against, so they are
for comparing schedulers, not for reading as measurements. No interference between co-located jobs is modelled.

%======================================================================
\section{Next step}
%======================================================================
Phase~6: DRF and Tetris, validated against Synergy (Fig.~13); real power inputs with sources (spec sheets and a
SPECpower result for the servers); and our own studies on the Philly trace with its real arrivals (Fig.~6) and on
newer traces, using the comparison runner and the decision report.

\end{document}
